\documentclass[11pt,a4paper]{article}

\usepackage[margin=2.25cm]{geometry}
\usepackage{amsmath,amssymb,mathtools,bm}
\usepackage{booktabs}
\usepackage{array}
\usepackage{microtype}
\usepackage[hidelinks]{hyperref}
\usepackage[nameinlink,noabbrev]{cleveref}

\allowdisplaybreaks
\setlength{\parindent}{0pt}
\setlength{\parskip}{0.55em}

\newcommand{\ii}{\mathrm{i}}
\newcommand{\e}{\mathrm{e}}
\newcommand{\Tr}{\operatorname{tr}}
\newcommand{\Ree}{\operatorname{Re}}
\newcommand{\Imm}{\operatorname{Im}}
\newcommand{\diag}{\operatorname{diag}}
\newcommand{\mat}[1]{\bm{#1}}
\newcommand{\sharpconj}[1]{#1^{\sharp}}
\newcommand{\fminus}{f_{1001}}
\newcommand{\fplus}{f_{1010}}

\title{Relation Between Mais--Ripken Lattice Functions and Linear-Coupling Resonance Driving Terms}
\author{}
\date{\today}

\begin{document}
\maketitle

\begin{abstract}
We collect, in one set of conventions, the description of stable linear transverse coupling by the Mais--Ripken (MR) lattice functions, by the linear normal-form matrix $W$, and by the coupling resonance driving terms (RDTs) $f_{1001}$ and $f_{1010}$.  The turn-by-turn orbit is written in both formalisms.  Position and momentum spectral-line ratios are derived, including the interference between the difference- and sum-resonance terms.  An exact algebraic correspondence is then given between the Edwards--Teng normalized coupling matrix and matrix-equivalent linear RDTs.  This correspondence yields the exact identities
\[
 \frac12\sqrt{\frac{\beta_{x2}}{\beta_{x1}}}
   =\left|f_{1001}-f_{1010}^{*}\right|,
 \qquad
 \frac12\sqrt{\frac{\beta_{y1}}{\beta_{y2}}}
   =\left|f_{1001}-f_{1010}\right|,
\]
and analogous relations for normalized momenta with the plus signs.  The perturbative RDTs obtained from a sum over lattice sources coincide with the matrix-equivalent quantities to first order in the coupling.  The derivation assumes stable, on-momentum, four-dimensional symplectic motion in canonical coordinates $(x,p_x,y,p_y)$ and adopts the Xsuite convention for complex Courant--Snyder variables and RDTs.
\end{abstract}

\section{Scope and conventions}

We restrict the discussion to the transverse canonical phase-space vector
\begin{equation}
 \mat z=(x,p_x,y,p_y)^T,
\end{equation}
with symplectic matrix
\begin{equation}
 \mat S=\diag(\mat J,\mat J),
 \qquad
 \mat J=\begin{pmatrix}0&1\\-1&0\end{pmatrix}.
\end{equation}
The one-turn map is assumed stable and symplectic.  Mode 1 is labeled horizontal-like and mode 2 vertical-like.  A superscript $*$ denotes complex conjugation.  For a real $2\times2$ matrix $A$, the symbol
\begin{equation}
 A^{\sharp}\equiv-\mat J A^T\mat J
\end{equation}
denotes its symplectic conjugate; it must not be confused with complex conjugation.

The exact matrix relations below apply to a purely transverse $4\times4$ symplectic problem.  Longitudinal coupling, radiation damping, non-symplectic maps, and the distinction between canonical momentum and geometric slope require additional terms.  The complex-variable and RDT signs follow the Xsuite physics-manual convention~\cite{XsuiteManual}.

\section{Turn-by-turn motion in the Mais--Ripken formalism}
\label{sec:mr}

\subsection{Linear normal form and the $W$ matrix}

At observation point $s$, let $M(s)$ be the one-turn map.  Its real linear normal form is written as
\begin{equation}
 M(s)=W(s)\,R(\mu_1,\mu_2)\,W^{-1}(s),
 \label{eq:normal-form}
\end{equation}
where
\begin{equation}
 R(\mu_1,\mu_2)=\diag\!\left(R_2(\mu_1),R_2(\mu_2)\right),
 \qquad
 R_2(\mu)=
 \begin{pmatrix}
  \cos\mu&\sin\mu\\
  -\sin\mu&\cos\mu
 \end{pmatrix},
\end{equation}
and
\begin{equation}
 W^T\mat S W=\mat S.
\end{equation}
The normal-mode phases at turn $N$ are
\begin{equation}
 \Phi_j(N)=2\pi Q_jN+\phi_{j0},
 \qquad j=1,2.
\end{equation}
Using actions $J_j$, the normalized mode coordinates are chosen as
\begin{equation}
 \widehat{\mat z}(N)=
 \begin{pmatrix}
  \sqrt{2J_1}\cos\Phi_1\\
  -\sqrt{2J_1}\sin\Phi_1\\
  \sqrt{2J_2}\cos\Phi_2\\
  -\sqrt{2J_2}\sin\Phi_2
 \end{pmatrix},
 \qquad
 \mat z(N)=W\widehat{\mat z}(N).
 \label{eq:mr-orbit-matrix}
\end{equation}

Define the complex eigenvectors
\begin{equation}
 \mat v_j=W_{:,2j-1}+\ii W_{:,2j}.
\end{equation}
Then the complete turn-by-turn orbit is
\begin{equation}
 \boxed{
 \mat z(N)=\sum_{j=1}^{2}\sqrt{2J_j}\,
 \Ree\!\left[\mat v_j\e^{\ii\Phi_j(N)}\right].
 }
 \label{eq:mr-orbit-complex}
\end{equation}
This form is independent of the phase gauge used for each eigenvector, provided the same gauge is used consistently in $W$ and in $\Phi_j$.

\subsection{MR lattice functions from $W$}

Partition $W$ into $2\times2$ blocks by physical plane and normal mode,
\begin{equation}
 W=
 \begin{pmatrix}
  W_{x1}&W_{x2}\\
  W_{y1}&W_{y2}
 \end{pmatrix}.
\end{equation}
For $q\in\{x,y\}$ and $j\in\{1,2\}$, define the projected mode matrix
\begin{equation}
 \Sigma_{qj}\equiv W_{qj}W_{qj}^T
 =
 \begin{pmatrix}
  \beta_{qj}&-\alpha_{qj}\\
  -\alpha_{qj}&\gamma_{qj}
 \end{pmatrix}.
 \label{eq:mr-sigma}
\end{equation}
Equivalently, if the two columns of $W_{qj}$ are denoted by $(w_{q,2j-1},w_{q,2j})$ and $(w_{p_q,2j-1},w_{p_q,2j})$, then
\begin{align}
 \beta_{qj}&=w_{q,2j-1}^2+w_{q,2j}^2,\\
 \alpha_{qj}&=-w_{q,2j-1}w_{p_q,2j-1}
               -w_{q,2j}w_{p_q,2j},\\
 \gamma_{qj}&=w_{p_q,2j-1}^2+w_{p_q,2j}^2.
 \label{eq:mr-functions-from-w}
\end{align}
Unlike an uncoupled Courant--Snyder triplet, an individual projected triplet need not satisfy $\beta\gamma-\alpha^2=1$.
These projected functions are the Mais--Ripken lattice functions~\cite{MaisRipken,LebedevBogacz}.

Writing
\begin{equation}
 v_{qj}=W_{q,2j-1}+\ii W_{q,2j}
       =\sqrt{\beta_{qj}}\e^{\ii\chi_{qj}},
\end{equation}
the position and canonical-momentum signals are
\begin{align}
 q(N)&=\sum_{j=1}^{2}\sqrt{2J_j\beta_{qj}}
       \cos\!\left(\Phi_j+\chi_{qj}\right),
 \label{eq:mr-position}\\
 p_q(N)&=\sum_{j=1}^{2}\sqrt{2J_j\gamma_{qj}}
       \cos\!\left(\Phi_j+\chi_{p_qj}\right),
 \label{eq:mr-momentum}
\end{align}
where $\chi_{p_qj}=\arg(W_{p_q,2j-1}+\ii W_{p_q,2j})$.  To avoid confusion
with the ET normalizing matrices $A_q$, let $\mathcal A_u(Q_j)$ denote the
amplitude of the Fourier line at tune $Q_j$ in a signal $u(N)$.  The position
and canonical-momentum line amplitudes are therefore
\begin{equation}
 \mathcal A_q(Q_j)=\sqrt{2J_j\beta_{qj}},
 \qquad
 \mathcal A_{p_q}(Q_j)=\sqrt{2J_j\gamma_{qj}}.
 \label{eq:mr-line-amplitudes}
\end{equation}

\section{Edwards--Teng decoupling and modal Twiss parameters}
\label{sec:et-twiss}

\subsection{Decoupled modal maps}

The Edwards--Teng construction decouples the physical four-dimensional map
into two stable two-dimensional modal maps.  We denote the horizontal-like
modal map by $M_x$ (mode 1) and the vertical-like modal map by $M_y$ (mode 2).
Here and below,
\begin{equation}
 I_2=\begin{pmatrix}1&0\\0&1\end{pmatrix}
\end{equation}
denotes the $2\times2$ identity matrix.  In the Edwards--Teng convention used
here~\cite{EdwardsTeng}, the decoupling matrix is
\begin{equation}
 T_{\mathrm{ET}}
 =\rho
 \begin{pmatrix}
  I_2&R_{\mathrm{ET}}^{\sharp}\\
  -R_{\mathrm{ET}}&I_2
 \end{pmatrix},
 \qquad
 \rho=\frac{1}{\sqrt{1+\det R_{\mathrm{ET}}}}>0.
 \label{eq:tet}
\end{equation}
Let $T_{\mathrm{ET}}$ map modal coordinates $\mat z_{\mathrm{ET}}$ to physical
coordinates $\mat z$:
\begin{equation}
 \mat z=T_{\mathrm{ET}}\mat z_{\mathrm{ET}}.
\end{equation}
It block-diagonalizes the physical one-turn map $M$ according to
\begin{equation}
 T_{\mathrm{ET}}^{-1}MT_{\mathrm{ET}}
 =\begin{pmatrix}M_x&0\\0&M_y\end{pmatrix},
 \label{eq:et-decoupled-map}
\end{equation}
or, equivalently,
\begin{equation}
 M=T_{\mathrm{ET}}
 \begin{pmatrix}M_x&0\\0&M_y\end{pmatrix}
 T_{\mathrm{ET}}^{-1}.
\end{equation}
Thus $M_x$ and $M_y$ are the diagonal blocks of the transformed, decoupled
map; in a coupled lattice they are not the physical diagonal blocks of $M$.

For $q\in\{x,y\}$, write
\begin{equation}
 M_q=A_qR_2(\mu_q)A_q^{-1},
 \label{eq:et-modal-map}
\end{equation}
where $R_2$ is the rotation matrix defined above and
\begin{equation}
 A_q=
 \begin{pmatrix}
  \sqrt{\beta_q}&0\\
  -\alpha_q/\sqrt{\beta_q}&1/\sqrt{\beta_q}
 \end{pmatrix}.
 \label{eq:Aq}
\end{equation}
The parameters $(\beta_x,\alpha_x)$ are therefore the Edwards--Teng Twiss
parameters of the horizontal-like mode, and $(\beta_y,\alpha_y)$ are those of
the vertical-like mode.  They are not the projected Mais--Ripken functions
$(\beta_{x1},\alpha_{x1})$ and $(\beta_{y2},\alpha_{y2})$.

Its inverse is
\begin{equation}
 A_q^{-1}=
 \begin{pmatrix}
  \beta_q^{-1/2}&0\\
  \alpha_q\beta_q^{-1/2}&\beta_q^{1/2}
 \end{pmatrix}
 \label{eq:Aq-inverse}
\end{equation}
and maps physical modal coordinates in plane $q$ to Edwards--Teng normalized
coordinates:
\begin{equation}
 \begin{pmatrix}\widetilde q\\\widetilde p_q\end{pmatrix}
 =A_q^{-1}\begin{pmatrix}q\\p_q\end{pmatrix}.
 \label{eq:cs-normalization}
\end{equation}
In the uncoupled limit, these modal Twiss parameters reduce to the usual
horizontal and vertical Courant--Snyder Twiss parameters.

\subsection{Relation between $W$ and the Edwards--Teng transformation}

The matrix $A_q^{-1}$ maps modal coordinates to normalized coordinates,
whereas $A_q$ performs the inverse operation.  At this point it is important
not to assume that the normalized coordinates associated with $W$ and with
the ET construction have the same phase gauge.  Denote them separately by
\begin{equation}
 \widehat{\mat z}_W=W^{-1}\mat z,
 \qquad
 \widehat{\mat z}_{\mathrm{ET}}
 =\diag(A_x^{-1},A_y^{-1})T_{\mathrm{ET}}^{-1}\mat z.
 \label{eq:two-normalized-coordinates}
\end{equation}
Define
\begin{equation}
 \mathcal A\equiv\diag(A_x,A_y),
 \qquad
 U\equiv T_{\mathrm{ET}}\mathcal A.
 \label{eq:U-definition}
\end{equation}
The second definition in \cref{eq:two-normalized-coordinates} is then
equivalent to
\begin{equation}
 \mat z=U\widehat{\mat z}_{\mathrm{ET}}.
 \label{eq:z-from-et-normalized}
\end{equation}

Both $U$ and $W$ reduce the physical map to the same block rotation.  To see
this explicitly, use $U^{-1}=\mathcal A^{-1}T_{\mathrm{ET}}^{-1}$ and substitute
the ET block-diagonalization from \cref{eq:et-decoupled-map}:
\begin{align}
 U^{-1}MU
 &=\mathcal A^{-1}T_{\mathrm{ET}}^{-1}
   MT_{\mathrm{ET}}\mathcal A
 \nonumber\\
 &=\mathcal A^{-1}\diag(M_x,M_y)\mathcal A
 \nonumber\\
 &=\diag(A_x^{-1}M_xA_x,A_y^{-1}M_yA_y)
 \nonumber\\
 &=\diag(R_2(\mu_1),R_2(\mu_2))
 \nonumber\\
 &=R(\mu_1,\mu_2).
 \label{eq:U-normalizes-M}
\end{align}
The penultimate equality follows directly from
$M_q=A_qR_2(\mu_q)A_q^{-1}$ in \cref{eq:et-modal-map}.
Likewise, \cref{eq:normal-form} gives
\begin{equation}
 W^{-1}MW=R(\mu_1,\mu_2).
\end{equation}
Consequently, the matrix
\begin{equation}
 D\equiv U^{-1}W
 \label{eq:D-definition}
\end{equation}
commutes with the block rotation:
\begin{equation}
 DR(\mu_1,\mu_2)=R(\mu_1,\mu_2)D.
\end{equation}
For distinct, nondegenerate tunes and a fixed ordering of the two modes, a
symplectic matrix with this property consists of an independent phase rotation
in each mode:
\begin{equation}
 D=\diag\!\left(R_2(\theta_1),R_2(\theta_2)\right).
 \label{eq:D-phase-rotations}
\end{equation}
It follows from $W=UD$ that the two normalized vectors are related by
\begin{equation}
 \widehat{\mat z}_{\mathrm{ET}}=D\widehat{\mat z}_W.
 \label{eq:normalized-coordinate-gauges}
\end{equation}
Thus they differ only by shifts of the two normal-mode phases.  These shifts
do not alter the actions or the physical orbit.

The columns of $W$ may now be phase-rotated so that
$D=\diag(I_2,I_2)$.  This choice is
the Edwards--Teng phase gauge.  In this gauge
\begin{equation}
 \widehat{\mat z}_{\mathrm{ET}}=\widehat{\mat z}_W
 \equiv\widehat{\mat z},
\end{equation}
and \cref{eq:U-definition} gives
\begin{equation}
 \boxed{
 W=T_{\mathrm{ET}}\diag(A_x,A_y).
 }
 \label{eq:w-from-tet}
\end{equation}
All equations below that involve the individual columns or blocks of $W$ use
this phase gauge.
Substituting the explicit ET matrix from \cref{eq:tet} into
\cref{eq:w-from-tet} first gives
\begin{equation}
 \begin{aligned}
 W
 &=\rho
 \begin{pmatrix}
  I_2&R_{\mathrm{ET}}^{\sharp}\\
  -R_{\mathrm{ET}}&I_2
 \end{pmatrix}
 \begin{pmatrix}A_x&0\\0&A_y\end{pmatrix}\\
 &=\begin{pmatrix}
  \rho A_x&\rho R_{\mathrm{ET}}^{\sharp}A_y\\
  -\rho R_{\mathrm{ET}}A_x&\rho A_y
 \end{pmatrix}.
 \end{aligned}
 \label{eq:w-et-block-product}
\end{equation}
To factor the modal normalizers $A_x$ and $A_y$ from the left, define the
ET-normalized coupling matrix
\begin{equation}
 \boxed{
 \mathcal C=\rho A_x^{-1}R_{\mathrm{ET}}^{\sharp}A_y.
 }
 \label{eq:c-from-ret}
\end{equation}
The upper-right block in \cref{eq:w-et-block-product} is consequently
\begin{equation}
 A_x\mathcal C=\rho R_{\mathrm{ET}}^{\sharp}A_y.
\end{equation}
For the lower-left block, use
\begin{equation}
 (PQ)^{\sharp}=Q^{\sharp}P^{\sharp},
 \qquad
 (R_{\mathrm{ET}}^{\sharp})^{\sharp}=R_{\mathrm{ET}},
 \qquad
 A_q^{\sharp}=A_q^{-1},
\end{equation}
where the last identity follows because $A_q$ is symplectic.  Taking the
symplectic conjugate of \cref{eq:c-from-ret} gives
\begin{align}
 \mathcal C^{\sharp}
 &=\rho
 \left(A_x^{-1}R_{\mathrm{ET}}^{\sharp}A_y\right)^{\sharp}
 \nonumber\\
 &=\rho A_y^{-1}R_{\mathrm{ET}}A_x,
\end{align}
and therefore
\begin{equation}
 A_y\mathcal C^{\sharp}=\rho R_{\mathrm{ET}}A_x.
\end{equation}
Using these two block identities, \cref{eq:w-et-block-product} becomes
\begin{equation}
 \boxed{
 W=
 \begin{pmatrix}A_x&0\\0&A_y\end{pmatrix}
 \begin{pmatrix}
  \rho I_2&\mathcal C\\
  -\mathcal C^{\sharp}&\rho I_2
 \end{pmatrix}.
 }
 \label{eq:w-factorization}
\end{equation}
Symplecticity requires
\begin{equation}
 \rho^2+\det\mathcal C=1.
 \label{eq:rho-c-condition}
\end{equation}
Equivalently, if $W$ is partitioned as in \cref{eq:w-factorization}, its
phase-gauged blocks obey
\begin{equation}
 \boxed{
 A_x^{-1}W_{x1}=\rho I_2,\qquad
 A_x^{-1}W_{x2}=\mathcal C,\qquad
 A_y^{-1}W_{y1}=-\mathcal C^{\sharp},\qquad
 A_y^{-1}W_{y2}=\rho I_2.
 }
 \label{eq:w-blocks-normalized}
\end{equation}
For later use, define
\begin{equation}
 u\equiv\det\mathcal C
 =\frac{\det R_{\mathrm{ET}}}{1+\det R_{\mathrm{ET}}}
 =1-\rho^2.
 \label{eq:u-def}
\end{equation}
For the ordinary trigonometric ET branch, $u=\sin^2\phi$ and
$\rho^2=\cos^2\phi$, where $\phi$ is the Teng rotation angle.

\subsection{Relation between Edwards--Teng and Mais--Ripken functions}

Using the definition $\Sigma_{qj}=W_{qj}W_{qj}^T$ from
\cref{eq:mr-sigma} and substituting the four blocks of $W$ from
\cref{eq:w-factorization} gives
\begin{align}
 \Sigma_{x1}&=\rho^2 A_xA_x^T,
 &\Sigma_{x2}&=A_x\mathcal C\mathcal C^TA_x^T,
 \label{eq:sigmax-factorization}\\
 \Sigma_{y1}&=A_y\mathcal C^{\sharp}(\mathcal C^{\sharp})^TA_y^T,
 &\Sigma_{y2}&=\rho^2 A_yA_y^T.
 \label{eq:sigmay-factorization}
\end{align}
With
\begin{equation}
 \gamma_x=\frac{1+\alpha_x^2}{\beta_x},
 \qquad
 \gamma_y=\frac{1+\alpha_y^2}{\beta_y},
\end{equation}
the principal MR functions are therefore
\begin{align}
 \beta_{x1}&=\rho^2\beta_x,
 &\alpha_{x1}&=\rho^2\alpha_x,
 &\gamma_{x1}&=\rho^2\gamma_x,
 \label{eq:mr-et-main-x}\\
 \beta_{y2}&=\rho^2\beta_y,
 &\alpha_{y2}&=\rho^2\alpha_y,
 &\gamma_{y2}&=\rho^2\gamma_y.
 \label{eq:mr-et-main-y}
\end{align}
The determinant of either principal projected block is therefore
\begin{equation}
 \beta_{x1}\gamma_{x1}-\alpha_{x1}^2
 =\beta_{y2}\gamma_{y2}-\alpha_{y2}^2
 =\rho^4.
 \label{eq:main-determinants}
\end{equation}
Consequently, $\rho$ is obtained directly from the principal MR functions:
\begin{equation}
 \boxed{
 \rho^2
 =\sqrt{\beta_{x1}\gamma_{x1}-\alpha_{x1}^2}
 =\sqrt{\beta_{y2}\gamma_{y2}-\alpha_{y2}^2}.
 }
 \label{eq:rho-from-mr}
\end{equation}
The ET modal Twiss parameters then follow by division by $\rho^2$:
\begin{equation}
 \boxed{
 \beta_x=\frac{\beta_{x1}}{\rho^2},\quad
 \alpha_x=\frac{\alpha_{x1}}{\rho^2},\qquad
 \beta_y=\frac{\beta_{y2}}{\rho^2},\quad
 \alpha_y=\frac{\alpha_{y2}}{\rho^2}.
 }
 \label{eq:et-from-main-mr}
\end{equation}

To make the remaining MR functions equally explicit, write
\begin{equation}
 \mathcal C=
 \begin{pmatrix}
  c_{11}&c_{12}\\c_{21}&c_{22}
 \end{pmatrix}.
 \label{eq:c-components}
\end{equation}
Then
\begin{align}
 \beta_{x2}
 &=\beta_x(c_{11}^2+c_{12}^2),
 \label{eq:betax2-et}\\
 \alpha_{x2}
 &=\alpha_x(c_{11}^2+c_{12}^2)
   -(c_{11}c_{21}+c_{12}c_{22}),
 \label{eq:alphax2-et}\\
 \gamma_{x2}
 &=\frac{
  (c_{21}-\alpha_xc_{11})^2
  +(c_{22}-\alpha_xc_{12})^2
 }{\beta_x},
 \label{eq:gammax2-et}
\end{align}
and, since
\begin{equation}
 -\mathcal C^{\sharp}
 =
 \begin{pmatrix}
  -c_{22}&c_{12}\\
  c_{21}&-c_{11}
 \end{pmatrix},
\end{equation}
one also has
\begin{align}
 \beta_{y1}
 &=\beta_y(c_{22}^2+c_{12}^2),
 \label{eq:betay1-et}\\
 \alpha_{y1}
 &=\alpha_y(c_{22}^2+c_{12}^2)
   +c_{21}c_{22}+c_{11}c_{12},
 \label{eq:alphay1-et}\\
 \gamma_{y1}
 &=\frac{
  (c_{21}+\alpha_yc_{22})^2
  +(c_{11}+\alpha_yc_{12})^2
 }{\beta_y}.
 \label{eq:gammay1-et}
\end{align}
The cross-plane projected determinants satisfy
\begin{equation}
 \begin{aligned}
 \beta_{x2}\gamma_{x2}-\alpha_{x2}^2
 &=\beta_{y1}\gamma_{y1}-\alpha_{y1}^2\\
 &=(1-\rho^2)^2.
 \end{aligned}
 \label{eq:cross-determinants}
\end{equation}

\section{Turn-by-turn motion in the RDT formalism}
\label{sec:rdt}

\subsection{Complex Courant--Snyder coordinates}

Using the Edwards--Teng normalized coordinates defined in
\cref{eq:cs-normalization}, the complex Courant--Snyder variables are
\begin{equation}
 h_{q,\pm}=\widetilde q\pm\ii\widetilde p_q
 =\sqrt{2I_q}\e^{\mp\ii\Phi_q}.
 \label{eq:complex-cs}
\end{equation}
Thus $h_{q,-}=\widetilde q-\ii\widetilde p_q$ and
\begin{equation}
 \widetilde q=\Ree(h_{q,-}),
 \qquad
 \widetilde p_q=-\Imm(h_{q,-}).
\end{equation}

The normal-form generating function is expanded as
\begin{equation}
 F=\sum_{jklm}f_{jklm}
 h_{x,+}^{j}h_{x,-}^{k}h_{y,+}^{l}h_{y,-}^{m}.
\end{equation}
To first order in the RDTs, the turn-by-turn motion is
\begin{align}
 h_{x,-}(s,N)
 &=\sqrt{2I_x}\e^{\ii\Phi_x}
 -2\ii\sum_{jklm}j f_{jklm}^{(s)}
 (2I_x)^{(j+k-1)/2}(2I_y)^{(l+m)/2}
 \nonumber\\[-0.2em]
 &\hspace{3.2cm}\times
 \e^{\ii[(1-j+k)\Phi_x+(m-l)\Phi_y]},
 \label{eq:general-hx-rdt}\\
 h_{y,-}(s,N)
 &=\sqrt{2I_y}\e^{\ii\Phi_y}
 -2\ii\sum_{jklm}l f_{jklm}^{(s)}
 (2I_x)^{(j+k)/2}(2I_y)^{(l+m-1)/2}
 \nonumber\\[-0.2em]
 &\hspace{3.2cm}\times
 \e^{\ii[(k-j)\Phi_x+(1-l+m)\Phi_y]}.
 \label{eq:general-hy-rdt}
\end{align}
Here $\Phi_x=2\pi Q_xN+\psi_{x,s,0}$ and similarly for $y$.

At first order in lattice perturbation theory, the local RDT at observation point $s$ can be written as
\begin{equation}
 f_{jklm}(s)=
 \frac{
  \displaystyle\sum_{w}h_{w,jklm}
  \e^{\ii[(j-k)\Delta\phi_{x,w\to s}+(l-m)\Delta\phi_{y,w\to s}]}
 }{
  1-\e^{2\pi\ii[(j-k)Q_x+(l-m)Q_y]}
 }.
 \label{eq:general-local-rdt}
\end{equation}
The coefficient $h_{w,jklm}$ is the corresponding Hamiltonian coefficient of source $w$.  For a thin skew quadrupole, $h_{w,1001}$ and $h_{w,1010}$ have the same local strength factor, proportional to
\[
 K_{1s,w}\sqrt{\beta_{x,w}\beta_{y,w}}/4;
\]
their different global behavior comes from their phase factors and resonance denominators~\cite{XsuiteManual}.

In particular,
\begin{align}
 f_{1001}(s)&=
 \frac{\displaystyle\sum_w h_{w,1001}
 \e^{\ii(\Delta\phi_{x,w\to s}-\Delta\phi_{y,w\to s})}}
 {1-\e^{2\pi\ii(Q_x-Q_y)}},
 \label{eq:f1001-source}\\
 f_{1010}(s)&=
 \frac{\displaystyle\sum_w h_{w,1010}
 \e^{\ii(\Delta\phi_{x,w\to s}+\Delta\phi_{y,w\to s})}}
 {1-\e^{2\pi\ii(Q_x+Q_y)}}.
 \label{eq:f1010-source}
\end{align}

\subsection{The four linear-coupling monomials}

Reality of the Hamiltonian implies
\begin{equation}
 f_{jklm}^{*}=f_{kjml}.
\end{equation}
Hence
\begin{equation}
 \boxed{
 f_{0110}=f_{1001}^{*},
 \qquad
 f_{0101}=f_{1010}^{*}.
 }
 \label{eq:rdt-conjugates}
\end{equation}
Only terms with $j>0$ enter $h_{x,-}$, whereas only terms with $l>0$ enter $h_{y,-}$.  With
\begin{equation}
 f\equiv f_{1001},
 \qquad
 g\equiv f_{1010},
\end{equation}
\cref{eq:general-hx-rdt,eq:general-hy-rdt} give
\begin{align}
 \boxed{
 h_{x,-}=\sqrt{2I_x}\e^{\ii\Phi_x}
 -2\ii\sqrt{2I_y}
 \left(f\e^{\ii\Phi_y}+g\e^{-\ii\Phi_y}\right)
 }
 \label{eq:hx-linear-coupling}\\
 \boxed{
 h_{y,-}=\sqrt{2I_y}\e^{\ii\Phi_y}
 -2\ii\sqrt{2I_x}
 \left(f^{*}\e^{\ii\Phi_x}+g\e^{-\ii\Phi_x}\right).
 }
 \label{eq:hy-linear-coupling}
\end{align}
The difference-resonance term $f_{1001}$ produces the opposite-plane line with the same sign of frequency; the sum-resonance term $f_{1010}$ produces the negative-frequency contribution.  In a real position signal these two contributions interfere at the same observed tune line.

\section{Observable position and momentum amplitudes}
\label{sec:observables}

\subsection{Normalized-coordinate signals}

Using $\widetilde q=\Ree h_{q,-}$ and $\widetilde p_q=-\Imm h_{q,-}$, the horizontal signal becomes
\begin{align}
 \widetilde x(N)
 &=\sqrt{2I_x}\cos\Phi_x
 +2\sqrt{2I_y}\,
 \Imm\!\left[(f-g^{*})\e^{\ii\Phi_y}\right],
 \label{eq:xtilde-rdt}\\
 \widetilde p_x(N)
 &=-\sqrt{2I_x}\sin\Phi_x
 +2\sqrt{2I_y}\,
 \Ree\!\left[(f+g^{*})\e^{\ii\Phi_y}\right].
 \label{eq:pxtilde-rdt}
\end{align}
Similarly,
\begin{align}
 \widetilde y(N)
 &=\sqrt{2I_y}\cos\Phi_y
 +2\sqrt{2I_x}\,
 \Imm\!\left[(f^{*}-g^{*})\e^{\ii\Phi_x}\right],
 \label{eq:ytilde-rdt}\\
 \widetilde p_y(N)
 &=-\sqrt{2I_y}\sin\Phi_y
 +2\sqrt{2I_x}\,
 \Ree\!\left[(f^{*}+g^{*})\e^{\ii\Phi_x}\right].
 \label{eq:pytilde-rdt}
\end{align}
The four basic combinations are summarized in \cref{tab:rdt-combinations}.

\begin{table}[htbp]
\centering
\caption{RDT combinations controlling the cross-plane mode projection.}
\label{tab:rdt-combinations}
\begin{tabular}{@{}ll@{}}
\toprule
Observable & Cross-mode complex coefficient\\
\midrule
$x$ from mode 2 & $f-g^{*}$\\
$\widetilde p_x$ from mode 2 & $f+g^{*}$\\
$y$ from mode 1 & $(f-g)^{*}$\\
$\widetilde p_y$ from mode 1 & $(f+g)^{*}$\\
\bottomrule
\end{tabular}
\end{table}

For a particle carrying both mode actions, the measured position-line ratios are therefore
\begin{align}
 \frac{\mathcal A_x(Q_2)}{\mathcal A_x(Q_1)}
 &=2\sqrt{\frac{I_y}{I_x}}\,|f-g^{*}|,
 \label{eq:measured-x-rdt-ratio}\\
 \frac{\mathcal A_y(Q_1)}{\mathcal A_y(Q_2)}
 &=2\sqrt{\frac{I_x}{I_y}}\,|f-g|.
 \label{eq:measured-y-rdt-ratio}
\end{align}
The corresponding normalized-momentum ratios are
\begin{align}
 \frac{\mathcal A_{\widetilde p_x}(Q_2)}{\mathcal A_{\widetilde p_x}(Q_1)}
 &=2\sqrt{\frac{I_y}{I_x}}\,|f+g^{*}|,
 \label{eq:measured-px-rdt-ratio}\\
 \frac{\mathcal A_{\widetilde p_y}(Q_1)}{\mathcal A_{\widetilde p_y}(Q_2)}
 &=2\sqrt{\frac{I_x}{I_y}}\,|f+g|.
 \label{eq:measured-py-rdt-ratio}
\end{align}

\subsection{MR--RDT amplitude identities}

From \cref{eq:mr-line-amplitudes},
\begin{align}
 \frac{\mathcal A_x(Q_2)}{\mathcal A_x(Q_1)}
 &=\sqrt{\frac{J_2}{J_1}}
   \sqrt{\frac{\beta_{x2}}{\beta_{x1}}},\\
 \frac{\mathcal A_y(Q_1)}{\mathcal A_y(Q_2)}
 &=\sqrt{\frac{J_1}{J_2}}
   \sqrt{\frac{\beta_{y1}}{\beta_{y2}}}.
\end{align}
Identifying $J_1=I_x$ and $J_2=I_y$ to the working order gives
\begin{align}
 \boxed{
 r_x\equiv\frac12\sqrt{\frac{\beta_{x2}}{\beta_{x1}}}
 =|f_{1001}-f_{1010}^{*}|
 }
 \label{eq:rx-main}\\
 \boxed{
 r_y\equiv\frac12\sqrt{\frac{\beta_{y1}}{\beta_{y2}}}
 =|f_{1001}-f_{1010}|
 .}
 \label{eq:ry-main}
\end{align}
These are exact identities for the matrix-equivalent RDTs defined in \cref{sec:matrix-rdt}; for the source-summed perturbative RDTs they hold to first order in the coupling.

The normalized momentum in plane $q$ is
\begin{equation}
 \widetilde p_q=\frac{\alpha_q q+\beta_q p_q}{\sqrt{\beta_q}}.
\end{equation}
Using the MR matrices, define
\begin{align}
 r_{\widetilde p_x}
 &\equiv\frac12
 \sqrt{
  \frac{
   \alpha_x^2\beta_{x2}
   -2\alpha_x\beta_x\alpha_{x2}
   +\beta_x^2\gamma_{x2}
  }{\beta_{x1}}
 },
 \label{eq:rpx-mr}\\
 r_{\widetilde p_y}
 &\equiv\frac12
 \sqrt{
  \frac{
   \alpha_y^2\beta_{y1}
   -2\alpha_y\beta_y\alpha_{y1}
   +\beta_y^2\gamma_{y1}
  }{\beta_{y2}}
 }.
 \label{eq:rpy-mr}
\end{align}
Then
\begin{align}
 \boxed{r_{\widetilde p_x}=|f_{1001}+f_{1010}^{*}|,}
 \label{eq:rpx-rdt}\\
 \boxed{r_{\widetilde p_y}=|f_{1001}+f_{1010}|.}
 \label{eq:rpy-rdt}
\end{align}
Thus the position projection selects the difference of the two complex coupling terms, while the normalized-momentum projection selects their sum.

The interference structure is explicit:
\begin{align}
 r_x^2&=|f|^2+|g|^2-2\Ree(fg),\\
 r_{\widetilde p_x}^2&=|f|^2+|g|^2+2\Ree(fg),\\
 r_y^2&=|f|^2+|g|^2-2\Ree(fg^{*}),\\
 r_{\widetilde p_y}^2&=|f|^2+|g|^2+2\Ree(fg^{*}).
 \label{eq:interference-identities}
\end{align}
Consequently, even a relatively small $f_{1010}$ produces a first-order modulation of an MR position ratio around $|f_{1001}|$.

\subsection{Physical canonical-momentum ratios}

The MR functions give directly
\begin{align}
 \frac{\mathcal A_{p_x}(Q_2)}{\mathcal A_{p_x}(Q_1)}
 &=\sqrt{\frac{J_2}{J_1}}
   \sqrt{\frac{\gamma_{x2}}{\gamma_{x1}}},\\
 \frac{\mathcal A_{p_y}(Q_1)}{\mathcal A_{p_y}(Q_2)}
 &=\sqrt{\frac{J_1}{J_2}}
   \sqrt{\frac{\gamma_{y1}}{\gamma_{y2}}}.
\end{align}
Because $p_q=(\widetilde p_q-\alpha_q\widetilde q)/\sqrt{\beta_q}$, the pure RDT combinations are mixed by $\alpha_q$.  The corresponding lattice ratios are
\begin{align}
 \boxed{
 \frac12\sqrt{\frac{\gamma_{x2}}{\gamma_{x1}}}
 =\frac{
  \left|f+g^{*}+\ii\alpha_x(f-g^{*})\right|
 }{\sqrt{1+\alpha_x^2}}
 }
 \label{eq:physical-px-ratio}\\
 \boxed{
 \frac12\sqrt{\frac{\gamma_{y1}}{\gamma_{y2}}}
 =\frac{
  \left|f+g-\ii\alpha_y(f-g)\right|
 }{\sqrt{1+\alpha_y^2}}
 .}
 \label{eq:physical-py-ratio}
\end{align}
At an ET waist, where $\alpha_q=0$, the physical and normalized momentum ratios coincide.

\section{Exact matrix-equivalent RDTs from $W$}
\label{sec:matrix-rdt}

\subsection{Definition from the normalized coupling matrix}

The Edwards--Teng construction in \cref{sec:et-twiss} provides $\rho$ and the
normalized coupling matrix $\mathcal C$ directly from $W$, before any RDT is
introduced.  In the Xsuite RDT convention, the corresponding exact
matrix-equivalent linear coupling RDTs are
\begin{align}
 \boxed{
 f_{1001}=\frac{
  c_{21}-c_{12}+\ii(c_{11}+c_{22})
 }{4\rho}
 }
 \label{eq:f1001-from-c}\\
 \boxed{
 f_{1010}=\frac{
  c_{12}+c_{21}+\ii(c_{11}-c_{22})
 }{4\rho}.
 }
 \label{eq:f1010-from-c}
\end{align}
Relative to the convention used in the original matrix--Hamiltonian formulas of Calaga, Tom\'as, and Franchi~\cite{CalagaTomasFranchi}, the Xsuite convention applies an overall minus complex conjugation to each of these two RDTs.

Conversely, with $f=f_{1001}$ and $g=f_{1010}$,
\begin{equation}
 \boxed{
 \mathcal C=2\rho
 \begin{pmatrix}
  \Imm(f+g)&\Ree(g-f)\\
  \Ree(f+g)&\Imm(f-g)
 \end{pmatrix}.
 }
 \label{eq:c-from-rdts}
\end{equation}
Its complex row combinations are
\begin{align}
 c_{11}+\ii c_{12}&=-2\ii\rho(f-g^{*}),
 \label{eq:c-row1-complex}\\
 c_{21}+\ii c_{22}&= 2\rho(f+g^{*}),
 \label{eq:c-row2-complex}\\
 (-\mathcal C^{\sharp})_{11}
 +\ii(-\mathcal C^{\sharp})_{12}
 &=-2\ii\rho(f^{*}-g^{*}),
 \label{eq:csharp-row1-complex}\\
 (-\mathcal C^{\sharp})_{21}
 +\ii(-\mathcal C^{\sharp})_{22}
 &=2\rho(f^{*}+g^{*}).
 \label{eq:csharp-row2-complex}
\end{align}
These four identities are the matrix counterpart of \cref{eq:xtilde-rdt,eq:pxtilde-rdt,eq:ytilde-rdt,eq:pytilde-rdt}.

Equation~\eqref{eq:rho-c-condition} becomes
\begin{equation}
 \det\mathcal C=4\rho^2\left(|f|^2-|g|^2\right),
 \qquad
 \boxed{
 \rho^2=\frac{1}{1+4(|f|^2-|g|^2)}.
 }
 \label{eq:rho-from-rdts}
\end{equation}
Thus $\rho=1+\mathcal O(f^2,g^2)$, explaining why it is absent in the first-order orbit formulas while canceling exactly from the amplitude ratios.
Combining \cref{eq:sigmax-factorization,eq:sigmay-factorization} with
\cref{eq:c-row1-complex,eq:c-row2-complex,eq:csharp-row1-complex,eq:csharp-row2-complex}
immediately yields the exact amplitude identities in
\cref{eq:rx-main,eq:ry-main,eq:rpx-rdt,eq:rpy-rdt}.

\subsection{Direct extraction of the coupling RDTs from the normal-form matrix}
\label{sec:direct-rdts-from-w}

The matrix-equivalent RDTs can be obtained directly from $W$ without first
reconstructing all four real entries of $\mathcal C$.  Assume the coordinate
order $(x,p_x,y,p_y)$ and that columns 3 and 4 of $W$ belong to the
vertical-like mode 2.

First form the complex mode-2 eigenvector
\begin{equation}
 \mat v_2^{(0)}=W_{:,3}+\ii W_{:,4}.
\end{equation}
Fix its phase by requiring its physical $y$ component to be real and positive:
\begin{equation}
 \mat v_2
 =\mat v_2^{(0)}
  \exp\!\left[-\ii\arg\!\left(v_{2,y}^{(0)}\right)\right],
 \qquad v_{2,y}>0.
 \label{eq:v2-et-phase-gauge}
\end{equation}
This selects the mode-2 part of the ET phase gauge discussed in
\cref{eq:D-phase-rotations}.  The mode-1 phase does not enter the following
calculation because its MR functions are phase invariant.

Next obtain $\rho$ from the principal horizontal MR functions and recover the
horizontal ET Twiss parameters:
\begin{align}
 \rho^2
 &=\sqrt{\beta_{x1}\gamma_{x1}-\alpha_{x1}^2},
 &\rho&=\sqrt{\rho^2},
 \label{eq:direct-rho-from-w}\\
 \beta_x&=\frac{\beta_{x1}}{\rho^2},
 &\alpha_x&=\frac{\alpha_{x1}}{\rho^2}.
 \label{eq:direct-et-twiss-from-w}
\end{align}
The quantities $\beta_{x1}$, $\alpha_{x1}$, and $\gamma_{x1}$ are obtained
from the first two columns of $W$ through
\cref{eq:mr-functions-from-w}.

Apply the inverse horizontal ET normalizer to the horizontal components of
$\mat v_2$:
\begin{equation}
 \begin{pmatrix}u_x\\u_{p_x}\end{pmatrix}
 =A_x^{-1}
 \begin{pmatrix}v_{2,x}\\v_{2,p_x}\end{pmatrix}
 =\begin{pmatrix}
  v_{2,x}/\sqrt{\beta_x}\\[0.3em]
  (\alpha_xv_{2,x}+\beta_xv_{2,p_x})/\sqrt{\beta_x}
 \end{pmatrix}.
 \label{eq:direct-normalized-v2}
\end{equation}
By \cref{eq:w-blocks-normalized}, $A_x^{-1}W_{x2}=\mathcal C$, so these two
complex numbers are
\begin{equation}
 u_x=c_{11}+\ii c_{12},
 \qquad
 u_{p_x}=c_{21}+\ii c_{22}.
 \label{eq:direct-c-rows}
\end{equation}
Substitution into \cref{eq:f1001-from-c,eq:f1010-from-c} finally gives
\begin{align}
 \boxed{
 f_{1001}=\frac{u_{p_x}+\ii u_x}{4\rho}
 }
 \label{eq:direct-f1001-from-w}\\
 \boxed{
 f_{1010}=\frac{(u_{p_x}-\ii u_x)^{*}}{4\rho}.
 }
 \label{eq:direct-f1010-from-w}
\end{align}
This is the direct route used by an implementation based on the complex
mode-2 eigenvector of $W$.

\subsection{Perturbative and matrix-equivalent RDTs}

The source sum in \cref{eq:general-local-rdt} is a first-order perturbative quantity.  The definitions in \cref{eq:f1001-from-c,eq:f1010-from-c} are exact algebraic coordinates of the linear symplectic coupling matrix.  In weak coupling,
\begin{equation}
 \rho=1+\mathcal O(K_{1s}^2),
 \qquad
 \mathcal C=\mathcal O(K_{1s}),
\end{equation}
so the two definitions agree at first order.  At stronger coupling, or very close to resonance, one should state explicitly whether $f_{1001}$ and $f_{1010}$ denote source-summed perturbative RDTs or matrix-equivalent RDTs extracted from the exact one-turn map.

\section{Conclusions}

The MR and RDT descriptions encode the same linear coupled motion, but emphasize different objects.  The MR functions are quadratic projections of the exact normal-mode eigenvectors and directly give physical position and momentum amplitudes.  The RDTs separate the difference- and sum-resonance content and expose their phases.  Their observable combinations are not $|f_{1001}|$ alone:
\begin{equation}
 \begin{array}{c|c}
 \text{projection}&\text{controlling complex combination}\\ \hline
 x\leftarrow\text{mode 2}&f_{1001}-f_{1010}^{*}\\
 \widetilde p_x\leftarrow\text{mode 2}&f_{1001}+f_{1010}^{*}\\
 y\leftarrow\text{mode 1}&f_{1001}-f_{1010}\\
 \widetilde p_y\leftarrow\text{mode 1}&f_{1001}+f_{1010}
 \end{array}
\end{equation}
This explains both the phase-dependent modulation of MR cross-plane beta ratios and the possible appearance of coupling in momentum while the corresponding position projection remains small.

\begin{thebibliography}{99}

\bibitem{MaisRipken}
H.~Mais and G.~Ripken,
\emph{Theory of Coupled Synchro-Betatron Oscillations I},
DESY Report DESY-M-82-05 (1982).

\bibitem{EdwardsTeng}
D.~A.~Edwards and L.~C.~Teng,
``Parametrization of linear coupled motion in periodic systems,''
\emph{IEEE Trans. Nucl. Sci.} \textbf{20}, 885--888 (1973),
\href{https://doi.org/10.1109/TNS.1973.4327279}{doi:10.1109/TNS.1973.4327279}.

\bibitem{LebedevBogacz}
V.~A.~Lebedev and S.~A.~Bogacz,
``Betatron motion with coupling of horizontal and vertical degrees of freedom,''
\emph{JINST} \textbf{5}, P10010 (2010),
\href{https://doi.org/10.1088/1748-0221/5/10/P10010}{doi:10.1088/1748-0221/5/10/P10010},
\href{https://arxiv.org/abs/1207.5526}{arXiv:1207.5526}.

\bibitem{CalagaTomasFranchi}
R.~Calaga, R.~Tom\'as, and A.~Franchi,
``Betatron coupling: Merging Hamiltonian and matrix approaches,''
\emph{Phys. Rev. ST Accel. Beams} \textbf{8}, 034001 (2005),
\href{https://doi.org/10.1103/PhysRevSTAB.8.034001}{doi:10.1103/PhysRevSTAB.8.034001}.

\bibitem{XsuiteManual}
Xsuite Collaboration,
\emph{Xsuite Physics Manual}, Chapters 3--4,
\url{https://xsuite.github.io/xsuite/docs/physics_manual/physics_man.pdf}.

\end{thebibliography}

\end{document}
